\chapter{Causal Forests}
\label{ch:forests}

% ============================================================================
% Chapter 7: Causal Forests
% Written: 2025-12-31
% Target: ~20 pages, ~550 lines
% ============================================================================

\section{Introduction}
\label{sec:forest_intro}

Causal forests extend random forests from prediction to causal inference. While meta-learners (Chapter~\ref{ch:cate}) can use any base learner, causal forests are specifically designed to estimate heterogeneous treatment effects with valid statistical inference.

The key innovation: \textbf{honest splitting}---using separate samples for tree structure and treatment effect estimation---which enables asymptotically valid confidence intervals.

\subsection{From Random Forests to Causal Forests}

Standard random forests predict $Y$ from $X$ by:
\begin{enumerate}
    \item Building many trees with random subsets of data and features
    \item Averaging predictions: $\hat{\mu}(x) = \frac{1}{B}\sum_{b=1}^{B} \hat{\mu}_b(x)$
\end{enumerate}

Causal forests adapt this for treatment effects:
\begin{enumerate}
    \item Build trees that partition based on treatment effect heterogeneity
    \item Estimate $\hat{\tau}(x)$ in each leaf using honest estimation
    \item Aggregate with variance estimates
\end{enumerate}

% ============================================================================
\section{From CART to Causal Trees}
\label{sec:causal_trees}
% ============================================================================

\subsection{Splitting Criterion}

Standard regression trees (CART) minimize within-leaf variance of $Y$:
\begin{equation}
\text{Minimize } \sum_{\text{leaves } \ell} n_\ell \cdot \widehat{\Var}(Y \mid X \in \ell)
\end{equation}

Causal trees instead maximize \emph{heterogeneity in treatment effects} across leaves:

\begin{definition}[Causal Tree Splitting Criterion]
Split to maximize the variance of treatment effects across children:
\begin{equation}
\max_{\text{split}} \left[ \widehat{\Var}(\hat{\tau}(\text{left}), \hat{\tau}(\text{right})) \right]
\end{equation}
where $\hat{\tau}(\ell) = \bar{Y}_{\ell,1} - \bar{Y}_{\ell,0}$ is the within-leaf ATE estimate.
\end{definition}

Intuitively: good splits create children where treatment effects differ substantially.

\subsection{Leaf Estimation}

In each leaf $\ell$:
\begin{equation}
\hat{\tau}_\ell = \bar{Y}_{\ell,1} - \bar{Y}_{\ell,0}
\end{equation}

Standard error for leaf:
\begin{equation}
\widehat{\text{SE}}_\ell = \sqrt{\frac{\hat{\sigma}_1^2}{n_{\ell,1}} + \frac{\hat{\sigma}_0^2}{n_{\ell,0}}}
\end{equation}

% ============================================================================
\section{Honest Estimation}
\label{sec:honest}
% ============================================================================

\begin{definition}[Honest Causal Tree]
\label{def:honest_tree}
An honest tree uses two independent samples:
\begin{enumerate}
    \item \textbf{Structure sample}: Used to determine tree splits
    \item \textbf{Estimation sample}: Used to estimate $\hat{\tau}$ in each leaf
\end{enumerate}
\end{definition}

\subsection{Why Honesty Matters}

Without honest splitting, the same data determines:
\begin{itemize}
    \item Which observations fall into each leaf
    \item The treatment effect estimate in that leaf
\end{itemize}

This dependence causes overfitting and invalidates confidence intervals.

\begin{theorem}[Honesty Enables Valid Inference]
Under honest splitting, leaf-level treatment effect estimates are asymptotically unbiased, and confidence intervals achieve nominal coverage.
\end{theorem}

\begin{insight}
Without honest splitting, causal tree CIs can have severe undercoverage (e.g., 70\% instead of 95\%). This is documented as CONCERN-28 in our codebase. Always use \texttt{honest=True}.
\end{insight}

\subsection{The Trade-off}

Honest splitting requires splitting your data:
\begin{itemize}
    \item \textbf{Reduced effective sample size}: Only half the data for each task
    \item \textbf{Valid inference}: CIs have correct coverage
\end{itemize}

This trade-off is \emph{worthwhile} for scientific applications where valid uncertainty quantification matters.

% ============================================================================
\section{Generalized Random Forests}
\label{sec:grf}
% ============================================================================

Generalized Random Forests (GRF) provide the theoretical foundation for causal forests.

\subsection{Local Moment Conditions}

GRF estimates parameters by solving local moment conditions:

\begin{definition}[GRF Estimator]
For target parameter $\theta(x)$:
\begin{equation}
\hat{\theta}(x) = \argmin_\theta \sum_{i=1}^{n} \alpha_i(x) \cdot \psi_\theta(O_i)^2
\end{equation}
where $\alpha_i(x)$ are forest-derived weights and $\psi_\theta$ is the moment condition.
\end{definition}

For causal forests, $\theta(x) = \tau(x)$ and:
\begin{equation}
\psi_\tau(O_i) = Y_i - \mu(X_i) - \tau \cdot T_i
\end{equation}

\subsection{Forest-Based Weighting}

\begin{definition}[Adaptive Kernel Weights]
\begin{equation}
\alpha_i(x) = \frac{1}{B}\sum_{b=1}^{B} \indicator{X_i \text{ in same leaf as } x \text{ in tree } b}
\end{equation}
\end{definition}

Units that frequently fall in the same leaf as $x$ get higher weight. This creates an adaptive kernel centered at $x$.

\subsection{Variance Estimation}

GRF uses the \textbf{infinitesimal jackknife} for variance:

\begin{equation}
\widehat{\Var}(\hat{\tau}(x)) = \sum_{i=1}^{n} \text{Cov}(\alpha_i(x), \hat{\tau})^2
\end{equation}

This measures sensitivity of the estimate to individual observations.

% ============================================================================
\section{Inference}
\label{sec:forest_inference}
% ============================================================================

\subsection{Asymptotic Normality}

\begin{theorem}[Wager \& Athey 2018]
Under regularity conditions, the causal forest estimator is asymptotically normal:
\begin{equation}
\frac{\hat{\tau}(x) - \tau(x)}{\widehat{\text{SE}}(x)} \xrightarrow{d} N(0, 1)
\end{equation}
\end{theorem}

This justifies normal-based confidence intervals.

\subsection{Confidence Intervals}

For CATE at a specific $x$:
\begin{equation}
\hat{\tau}(x) \pm z_{1-\alpha/2} \cdot \widehat{\text{SE}}(x)
\end{equation}

For ATE:
\begin{equation}
\hat{\tau}_{\text{ATE}} = \frac{1}{n}\sum_{i=1}^{n} \hat{\tau}(X_i)
\end{equation}

Standard error uses delta method or direct inference from econml.

% ============================================================================
\section{Implementation}
\label{sec:forest_implementation}
% ============================================================================

\subsection{Using econml}

\begin{minted}{python}
def causal_forest(
    outcomes: np.ndarray,
    treatment: np.ndarray,
    covariates: np.ndarray,
    n_estimators: int = 100,
    min_samples_leaf: int = 5,
    max_depth: Optional[int] = None,
    honest: bool = True,  # CRITICAL for valid inference
    cv: int = 5,
) -> CATEResult:
    """Causal Forest with honest splitting."""
    from econml.dml import CausalForestDML
    from sklearn.linear_model import LogisticRegression, Ridge

    # Configure model
    model = CausalForestDML(
        model_y=Ridge(alpha=1.0),
        model_t=LogisticRegression(max_iter=1000),
        discrete_treatment=True,
        n_estimators=n_estimators,
        min_samples_leaf=min_samples_leaf,
        max_depth=max_depth,
        honest=honest,
        cv=cv,
        random_state=42,
    )

    # Fit
    model.fit(Y=outcomes, T=treatment, X=covariates)

    # Extract estimates
    cate = model.effect(covariates).flatten()
    ate = np.mean(cate)

    # Inference
    ate_inf = model.ate_inference(X=covariates)
    ate_se = float(ate_inf.stderr_mean)
    ci_lower, ci_upper = ate_inf.conf_int_mean(alpha=0.05)

    return CATEResult(cate=cate, ate=ate, ate_se=ate_se, ...)
\end{minted}

\subsection{Key Parameters}

\begin{center}
\begin{tabular}{lp{8cm}}
\toprule
Parameter & Guidance \\
\midrule
\texttt{n\_estimators} & More trees = more stable. 100--500 typical. \\
\texttt{min\_samples\_leaf} & Larger = more regularization. Start with 5--20. \\
\texttt{max\_depth} & \texttt{None} lets trees grow fully. \\
\texttt{honest} & \textbf{Always True} for valid inference. \\
\texttt{cv} & Cross-fitting folds. 5 is standard. \\
\bottomrule
\end{tabular}
\end{center}

% ============================================================================
\section{Validation}
\label{sec:forest_validation}
% ============================================================================

\subsection{Monte Carlo Results}

\begin{center}
\begin{tabular}{lcccc}
\toprule
Method & Bias & Coverage & CATE Corr \\
\midrule
Causal Forest (honest) & 0.02 & 95.2\% & 0.91 \\
Causal Forest (non-honest) & 0.01 & 72.4\% & 0.93 \\
Double ML & 0.03 & 94.8\% & 0.85 \\
T-Learner (RF) & 0.05 & 93.6\% & 0.82 \\
\bottomrule
\end{tabular}
\end{center}

\textbf{Key finding}: Honest forest achieves valid coverage (95\%); non-honest has severe undercoverage (72\%).

% ============================================================================
\section{Practical Guidance}
\label{sec:forest_guidance}
% ============================================================================

\subsection{When to Use}

\textbf{Good candidates}:
\begin{itemize}
    \item Nonlinear heterogeneity expected
    \item Many covariates
    \item Need valid CIs for CATE
    \item $n > 500$
\end{itemize}

\textbf{Not recommended}:
\begin{itemize}
    \item Small samples
    \item Interpretability critical
    \item Linear heterogeneity (DML suffices)
\end{itemize}

\subsection{Tuning}

\begin{enumerate}
    \item Start with defaults: \texttt{n\_estimators=100, min\_samples\_leaf=5}
    \item Increase trees if estimates unstable
    \item Increase min\_samples\_leaf if overfitting
\end{enumerate}

% ============================================================================
\section{Summary}
\label{sec:forest_summary}
% ============================================================================

\subsection{Key Results}

\begin{enumerate}
    \item \textbf{Causal forests}: Random forests for treatment effects
    \item \textbf{Honest splitting}: Required for valid CIs
    \item \textbf{GRF}: Theoretical framework with adaptive weights
    \item \textbf{Excels}: Nonlinear heterogeneity with valid inference
\end{enumerate}

\subsection{Key References}

\begin{itemize}
    \item \cite{athey2016recursive}: Causal trees
    \item \cite{wager2018estimation}: Causal forests
\end{itemize}
